\section{Verwandte Arbeiten}
\label{sec:related-work}

HAWQ \citep{dong_hawq_2019} und HAWQ-V2 \citep{dong_hawqv2_2020} etablieren die Hessian-Krümmung als Signal fuer die Bit-Zuteilung in gemischt-präzisen Quantisierungsschemata (\S\ref{sec:fund-hawq}). Beide Arbeiten validieren ihren Ansatz End-to-End, das heißt, nach der Bit-Zuteilung wird die Genauigkeit des gesamten Netzes gemessen und nicht die Vorhersagegüte des Scores für einzelne Schichten. Diese Validierungsstrategie ist für den ursprünglichen Zweck, eine gute globale Zuteilung zu finden, angemessen. Sie lässt jedoch offen, ob der Score tatsächlich lokal korrekt ist, denn ein Netz kann trotz einer schlecht kalibrierten Einzelschicht-Rangfolge insgesamt gut abschneiden, wenn genügend Gesamt-Bit-Budget vorhanden ist oder sich die Fehler zwischen den Schichten teilweise kompensieren.
\par\medskip

\citet{hill2026saturationmakesquantizationerror} testen als erste Arbeit direkt, ob die HAWQ-V2-Metrik $\mathrm{Tr}(H)\cdot\|\Delta W\|^2$ mit tatsächlich gemessenem Quantisierungsschaden korreliert und nicht nur mit der End-to-End-Genauigkeit nach der Bit-Zuteilung. Das Setting dieser Arbeit ist jedoch ein großes Sprachmodell (LLM) unter uniformer Quantisierung und kein Vision-CNN unter PoT. Da PoT, wie in \S\ref{sec:fund-quant} erläutert, ein grundlegend anderes, nicht-uniformes Fehlerverhalten aufweist, ist unklar, ob sich die dortigen Befunde auf PoT-CNNs übertragen lassen. Diese Übertragung, und ihr Scheitern in einem entscheidenden Punkt (\S\ref{sec:discussion}), ist ein zentraler Beitrag dieses Projekts.
\par\medskip

\citet{miyashita_pot_2016} fuehren PoT als hardware-freundliche logarithmische Gewichtsrepraesentation ein, \citet{lee_lognet_2017} zeigen den zugehoerigen Energie-Vorteil logarithmischer Berechnung für CNNs, und \citet{przewlocka_rus_pot_2022} motivieren PoT zusaetzlich für Hardware mit niedriger Bit-Breite. Neuere Arbeiten wie \citet{kim_nonzerogrid_2023} und \citet{li_denseshift_2023} verfeinern die PoT-Gitterformulierung selbst, um bei sehr niedriger Bit-Breite höhere End-to-End-Genauigkeit zu erreichen, und \citet{elgenedy_pot_llm_2025} überträgt PoT auf Sprachmodell-Gewichte. Alle diese Arbeiten fokussieren, wie \citet{miyashita_pot_2016} und \citet{przewlocka_rus_pot_2022}, auf die Machbarkeit, End-to-End-Genauigkeit und den Hardware-Vorteil von PoT und nicht auf eine schichtweise Sensitivitätsanalyse. In der PoT-Praxis, auch in diesem Projekt (\S\ref{sec:methodology}), ist es zudem üblich, die erste Faltungsschicht (\texttt{conv1}) entweder von der Quantisierung auszunehmen oder ihr besondere Aufmerksamkeit zu widmen, da sie mit ihren wenigen Eingabekanälen (meist drei Farbkanäle) und dem dadurch kleinen Fan-in empirisch besonders empfindlich reagiert. Dieses Projekt liefert mit der Random-Init-Kontrolle (Phase~2) und der Fusion-Analyse (Phase~3.5) eine systematische Erklärung dafür, warum \texttt{conv1} in dieser Hinsicht eine Sonderrolle einnimmt.
\par\medskip

SmoothQuant \citep{xiao_smoothquant_2023} und AWQ \citep{lin_awq_2024} zeigen, dass Aktivierungsschaden primär durch kanalweise Ausreißer getrieben wird, also durch einen von der Gewichtskruemmung unabhängigen Mechanismus (\S\ref{sec:fund-activation}). BRECQ \citep{li_brecq_2021} adressiert PTQ durch eine blockweise Rekonstruktion, die Gewichts- und Aktivierungsfehler gemeinsam optimiert, ohne explizit die Hessian-Krümmung als Sensitivitätssignal zu nutzen. Diese Arbeiten motivieren, warum dieses Projekt Gewichts- und Aktivierungsschaden explizit trennt (Phase~1, \S\ref{sec:phase1}), statt beide unter einem gemeinsamen Score zu behandeln.
\par\medskip

Zusammenfassend lässt sich festhalten, dass vor diesem Projekt kein publizierter empirischer Test existiert, der die HAWQ-V2-Formulierung direkt gegen gemessenen, per-Schicht-isolierten Gewichtsschaden für Vision-CNNs unter PoT-Quantisierung prüft.